% ================================================================
\subsection{The Renormalization Group of the Gram Energy}
\label{sec:rg-flow}
% ================================================================

The Cathedral's \lean{Geometry/Renormalization/} subsystem (18~files,
$\sim$6{,}000 lines) formalizes the \textbf{renormalization group flow}
of the Gram energy $F(N) = (v^T G v - 1) \cdot \ln N$.

\subsubsection{The RG Equation}

Numerical analysis reveals a scaling equation for the running coupling
$F(s)$ where $s = \ln N$:
\[
  \frac{dF}{ds} = \beta(s) \approx -\frac{1.76}{s^{1.82}}
\]
The beta function $\beta(s) < 0$ is {\em asymptotically free}:
the coupling weakens at large $N$ (high energy). The fixed point
is $L_1 = -\gamma - \ln(4\pi) \approx -3.108$, corresponding to
the B\'aez-Duarte constant regime.

\begin{correspondence}[RG Flow as Asymptotic Freedom]
\label{corr:rg-flow}
The negative beta function $\beta(s) < 0$ is the number-theoretic
analogue of \emph{asymptotic freedom} in QCD. In QCD, the coupling
constant $\alpha_s$ decreases at high energies (short distances),
and quarks become effectively free inside hadrons. Here, the Gram
energy decreases at large $N$ (many modes), and the quadratic form
$v^T G v$ approaches 1 from above---the vacuum energy flows to its
ground state.

The fixed point $L_1$ is the \emph{infrared fixed point} of the
prime number quantum field theory.
\end{correspondence}

\subsubsection{The Coprime Sector Decomposition}

The coprime/noncoprime decomposition (\lean{CoprimeSector.lean}) partitions
the double sum $v^T G v = \sum_{j,k} v_j G(j,k) v_k$ into:
\[
  v^T G v = \underbrace{E_{\mathrm{coprime}}}_{
    \gcd(j,k) = 1\text{ fibers}} +
  \underbrace{E_{\mathrm{non}}}_{
    \gcd(j,k) > 1\text{ fibers}}
\]

The coprime sector dominates and provides the destructive interference
(overcancellation) that drives $v^T G v \leq 1$. The noncoprime sector
contributes positive but subleading corrections via the Selberg
sieve structure.

\begin{correspondence}[Coprime Fibers as Interaction Channels]
Each coprime pair $(j,k)$ with $\gcd(j,k) = 1$ is an independent
scattering channel. The M\"obius function weights these channels
with signs $\mu(j)\mu(k) \in \{+1, -1\}$, creating the destructive
interference that is the physical content of the Riemann Hypothesis.
The noncoprime pairs $(j,k)$ with $\gcd(j,k) = d > 1$ reduce to
coprime pairs via $G(j,k) = G(j/d, k/d)$ (gauge reduction,
\lean{GCDReduction.lean}), contributing at order $1/d^2$.
\end{correspondence}

\subsubsection{The Selberg Bridge}

The Selberg sieve connection (\lean{SelbergBridge.lean}) provides the
PNT-level control needed to bound the noncoprime sector. The key insight
is that the GCD anatomy of the quadratic form mirrors the
\emph{Selberg $\Lambda^2$ sieve}: the weights $\mu(n)/n$ that appear
in the M\"obius log-cutoff witness are the natural Selberg weights for
the squarefree sieve.

This explains why $\beta(s) < 0$: each new squarefree stratum
$\{n : \omega(n) = k\}$ contributes \textbf{negative} interference to the
Gram energy, and the Euler product relay $\sum 1/d$ ensures the supply
of cancelling strata never runs out. The overcancellation is not
accidental---it is built into the multiplicative structure of $\NN$.

\subsubsection{Architecture Summary}

\begin{center}
\small
\begin{tabular}{@{}llr@{}}
\toprule
\textbf{Module} & \textbf{Physics Role} & \textbf{Lines} \\
\midrule
\lean{RGFlow.lean} & Beta function, fixed point & 809 \\
\lean{CoprimeSector.lean} & Channel decomposition & 472 \\
\lean{SelbergBridge.lean} & Sieve connection & 631 \\
\lean{MassRenormalization.lean} & $K_1 + L_1$ splitting & 305 \\
\lean{DiagonalCancellation.lean} & Self-energy subtraction & 163 \\
\lean{FinalAssembly.lean} & Graduation pipeline & 182 \\
\lean{OvercancellationGraduation.lean} & Crown axiom wiring & 227 \\
\midrule
\textbf{Total (18 files)} & & \textbf{6,041} \\
\bottomrule
\end{tabular}
\end{center}

All 18 files build with \textbf{zero errors, zero warnings}.
The subsystem connects the Mass Renormalization Theorem
(\S\ref{sec:mass-renormalization}) to the Glass Box graduation
(\S\ref{sec:glass-box-graduation}), providing the complete
renormalization group infrastructure for the prime number quantum field.
